\chapter*{Introduzione}
\addcontentsline{toc}{section}{\textbf{Introduzione}}

Il 29 luglio 2025, un terremoto di magnitudo \(\mathrm{M}\,8.8\) ha colpito la penisola della Kamchatka, nella Russia orientale.
Terremoti di grande magnitudo (\(\mathrm{M}\ge7\)) sono eventi particolarmente rari e non prevedibili, spesso catastrofici, di grande interesse sia scientifico che sociale. Il rischio per la popolazione è tendenzialmente molto elevato, non solo per l'evento sismico ma anche per gli effetti che il terremoto può produrre, come gli tsunami. Tra i più disastrosi si ricorda il terremoto \(\mathrm{M}\,9.1\) avvenuto a Sumatra, in Indonesia, nel dicembre 2004, che ha provocato un gigantesco tsunami. Questo ha causato circa \qty{230000}{} vittime, raggiungendo addirittura le coste di Kenya e Somalia. Tra i più noti vi è il terremoto \(\mathrm{M}\,9.0\) avvenuto a Tohoku, in Giappone nel marzo del 2011. Questo, insieme allo tsunami che ha scatenato, ha causato circa \qty{20000}{} vittime e provocato l'incidente nucleare nella centrale di Fukushima.

Fortunatamente, per il recente terremoto in Kamchatka non si sono verificate conseguenze così disastrose. Il terremoto ha sì prodotto uno tsunami, ma di intensità inferiore a quella prevista e le misure di prevenzione e evacuazione predisposte sono state altamente efficaci, causando una sola vittima.

Lo tsunami non è stato l'unico evento che ha seguito il terremoto. Nei giorni successivi, infatti, sono state registrate le eruzioni dei vulcani Klyuchevskoy e Krasheninnikov. L'innesco di eruzioni vulcaniche da parte di un terremoto è un argomento di ricerca tutt'ora aperto e fonte di studio.
Sono stati condotti diversi studi e analisi statistiche sulla possibile correlazione tra terremoti e vulcanismo. Tra i più recenti, \citet{Sawi2018ShortTermEQVolcanism} hanno concluso che l'attività vulcanica che si sviluppa pochi giorno dopo il sisma non è strettamente legata ad esso, mentre si nota un aumento più significativo dell'attività dai due ai cinque anni successivi. 
Nonostante gran parte delle eruzioni non dipenda da attività sismica, studiare questo fenomeno è comunque di grande importanza per migliorare la gestione del rischio \citep{Seropian2021EarthquakesVolcanoes}.
Ci sono due cause principali che possono indurre attività vulcanica a seguito di un terremoto: le vibrazioni del suolo legate a sforzi dinamici e le deformazioni permanenti del terreno, associate a sforzi statici.
In questa tesi si studia il secondo caso applicato al terremoto del 29 luglio in Kamchatka, con particolare riferimento alle eruzioni dei vulcani Klyuchevskoy e Krasheninnikov.

Nel \textit{Capitolo 1} viene tratta la descrizione geometrica di una faglia e introdotta brevemente la teoria tettonica delle placche.
Il \textit{Capitolo 2} presenta una descrizione della regione della Kamchatka, dell'evento sismico e dei vulcani studiati.
Si trattano poi, nel \textit{Capitolo 3}, sia i modelli usati per simulare gli effetti del terremoto, basati su dislocazioni rettangolari, sia la risalita del magma all'interno di fratture idrauliche, che prendono il nome di dicchi. I risultati della simulazione vengono esposti nel \textit{Capitolo 4}, per poi trarre le conclusioni nel \textit{Capitolo 5}.

